\section*{NeurIPS Paper Checklist}

\begin{enumerate}

\item {\bf Claims}
    \item[] Question: Do the main claims made in the abstract and introduction accurately reflect the paper's contributions and scope?
    \item[] Answer: \answerYes{}
    \item[] Justification: The contributions are (1) the methodological diagnosis that full-image identity metrics are structurally confounded under hard-composited edits, with outside-mask SSIM above $0.999$ by construction and replicating across blepharoplasty, rhinoplasty, and rhytidectomy; (2) the SurgicalScore benchmark protocol (five components A--E with weighted composite) and matched HDA paired splits ($N{=}27/21/9$) released as infrastructure; (3) cross-method, cross-procedure characterization on the $N{=}211$ rhinoplasty cohort showing no method achieves positive identity-space gain over the unedited input proxy, with \method{} producing the smallest residual gap ($-0.048$, paired permutation $p < 10^{-4}$ against each baseline); (4) \method{} as a depth-conditioned FLUX-Fill reference pipeline serving as a pre-consult goal-visualization tool, not a patient-specific outcome predictor. All four are directly traceable to evaluation files (\texttt{evaluation/baselines/strict\_n211\_all\_methods.json}, \texttt{strict\_n211\_baselines/}, \texttt{cross\_proc\_baselines/}, \texttt{cohort\_expansion\_2026-05-04/}). No claim exceeds what the referenced data support.
    \item[] Guidelines: \answerYes{} if all abstract and introduction claims are supported by experimental results or theoretical proofs.

\item {\bf Limitations}
    \item[] Question: Does the paper discuss the limitations of the work performed by the authors?
    \item[] Answer: \answerYes{}
    \item[] Justification: Section~\ref{sec:discussion} contains a Limitations paragraph naming the headline scope constraint (\method{} is a pre-consult goal-visualization tool, not a patient-specific outcome predictor); detailed limitations are in Appendix~\ref{app:limitations_full}, including: inside-mask accuracy does not exceed the pre-operative input proxy for any evaluated method; the matched test split is small ($N{=}21$); single-seed baselines introduce asymmetric variance; 2D geometry cannot represent three-dimensional rhinoplasty reshaping; gate thresholds are heuristic and not clinically validated; the FLUX.1-dev license is non-commercial; no cosmetic panel rating study has been conducted; and Monk Skin Tone coverage is preliminary with five of ten tones absent.
    \item[] Guidelines: \answerYes{} if limitations are discussed with appropriate depth.

\item {\bf Theory assumptions and proofs}
    \item[] Question: For each theoretical result, does the paper provide the full set of assumptions and a complete (and correct) proof?
    \item[] Answer: \answerNA{}
    \item[] Justification: No formal theoretical results are claimed. The relevant guarantee is the directly measured outside-mask SSIM exceeding $0.999$ across all evaluated cohorts, a property of the hard-mask composite $C = M \odot G + (1-M) \odot I$. An empirical Lipschitz characterization of ArcFace under masked perturbations ($L_{p95} \leq 1.2 \times 10^{-5}$ on $51$ cases $\times\,5$ noise magnitudes $=255$ perturbation samples; Appendix~\ref{app:sensitivity}) is reported as a numerical bound, not as a theorem.
    \item[] Guidelines: \answerYes{} if each theorem has a complete proof with explicit assumptions.

\item {\bf Experimental result reproducibility}
    \item[] Question: Does the paper fully disclose all information needed to reproduce the main experimental results?
    \item[] Answer: \answerYes{}
    \item[] Justification: Section~\ref{sec:experiments} specifies model names (FLUX.1-dev, FLUX.1-Fill-dev, FLUX.1-Kontext-dev, Depth Anything V2 Small, InsightFace buffalo\_l), inference hyperparameters (guidance 3.5, strength 0.75, ControlNet scale 0.5, 20 steps, BF16, $512{\times}512$), hardware (single NVIDIA L40S, 48\,GB), inference time (20\,s per image), and random seeds (seed 42 for splits and bootstrap). DPO training hyperparameters are in Appendix~\ref{app:dpo_hparams}. The released code package contains all evaluation harnesses.
    \item[] Guidelines: \answerYes{} if sufficient information is given to reproduce the main experimental results.

\item {\bf Open access to data and code}
    \item[] Question: Does the paper provide open access to the data and code, with sufficient instructions to reproduce the main results?
    \item[] Answer: \answerYes{}
    \item[] Justification: The SurgicalScore evaluation package, the \method{} pipeline code, and the HDA matched test split manifests/indices ($N{=}27$ blepharoplasty / $N{=}21$ rhinoplasty / $N{=}9$ rhytidectomy case identifiers and per-case metric outputs) are released at the supplementary material (zipped code bundle submitted alongside this paper PDF). The HDA \emph{images} are not redistributed; they are obtained from the original authors~\citep{rathgeb2020plastic} under their research-use terms. The supplementary repository contains evaluation scripts, preset configurations, split manifests, and a README with step-by-step reproduction instructions.
    \item[] Guidelines: \answerYes{} if a link to a code and data repository is provided.

\item {\bf Experimental setting/details}
    \item[] Question: Does the paper specify all training and test details necessary to understand the results?
    \item[] Answer: \answerYes{}
    \item[] Justification: Section~\ref{sec:experiments} reports data splits (matched HDA $N{=}27$ blepharoplasty / $N{=}21$ rhinoplasty / $N{=}9$ rhytidectomy paired pre/post pairs; $N{=}211$ external rhinoplasty cohort filtered from 457-pair ASPS+PCA pool by MediaPipe + ArcFace at $0.65$), all model versions, inference hyperparameters, and per-procedure evaluation denominators. The asymmetric ArcFace detection denominator (16/21 InsightFace detection successes on the matched pool) is disclosed with the rationale. Bootstrap parameters ($B{=}10{,}000$, seed 42) are stated. Nothing primary is omitted.
    \item[] Guidelines: \answerYes{} if all training and test details are reported.

\item {\bf Experiment statistical significance}
    \item[] Question: Does the paper report error bars suitably and correctly defined?
    \item[] Answer: \answerYes{}
    \item[] Justification: The headline $N{=}211$ ArcFace results in Table~\ref{tab:strict_n211} include 95\% percentile bootstrap CIs ($B{=}10{,}000$) over per-case scores. Table~\ref{tab:paired_n205} reports paired sign-flip permutation tests against each baseline on the four-way detected intersection ($N{=}205$, $p < 10^{-4}$ for ICEdit, FLUX.1-Kontext-dev, and InstructPix2Pix). The $N{=}211$ SurgicalScore comparison (Tables~\ref{tab:ss_strict_n211},~\ref{tab:ss_paired_strict}) carries per-case bootstrap CIs and paired permutation $p$-values on the four-way detected intersection ($N{=}188$). Mask-crop LPIPS (Table~\ref{tab:mask_crop}) and component correlation (Table~\ref{tab:component_corr}) carry bootstrap CIs and discriminative-power statistics. Sensitivity sweeps (Tables~\ref{tab:floor_sweep},~\ref{tab:loo_components},~\ref{tab:empirical_lipschitz}) defend the heuristic thresholds. Single-seed legacy baselines (Table~\ref{tab:m9_dpo}) do not have per-case bootstrap; the $\dagger$ notation distinguishes seed-level from per-case bootstrap.
    \item[] Guidelines: \answerYes{} if error bars are defined and reported or a clear rationale is given for why not.

\item {\bf Experiments compute resources}
    \item[] Question: Does the paper provide sufficient information on compute resources?
    \item[] Answer: \answerYes{}
    \item[] Justification: Section~\ref{sec:experiments} states single NVIDIA L40S (48\,GB), BF16 precision, 20 seconds per image inference. Appendix~\ref{app:dpo_hparams} discloses that DPO training ran on a single L40S for 2000 steps. The full matched-pool evaluation ($N{=}21$, 5 seeds) required approximately 35 GPU-minutes at this rate.
    \item[] Guidelines: \answerYes{} if compute is reported or a clear rationale is given.

\item {\bf Code of ethics}
    \item[] Question: Does the research conform with the NeurIPS Code of Ethics?
    \item[] Answer: \answerYes{}
    \item[] Justification: The work uses retrospective facial images, which are biometric and potentially sensitive. The paper uses the publicly available HDA Plastic Surgery Database under its published research-use terms; no human subjects were recruited or enrolled. External validation images (ASPS gallery, private clinical archive) are accessed under authorized release terms, used only for analysis, and not redistributed. We do not redistribute private clinical archive images or external validation images as part of the released code/data package. Released assets contain code, evaluation harnesses, preset definitions, and split manifests/indices. Qualitative images shown in the paper are included only under applicable source terms and are not part of the released supplementary data. No clinical claims are made; all outputs are labeled simulation only. The generative pipeline is not deployed or made publicly accessible in a form that could enable harm. PCA images are used under authorization from the providing clinical archive; prospective clinical deployment or patient-facing use would require institutional review and consent procedures.
    \item[] Guidelines: \answerYes{} if the work conforms with the NeurIPS Code of Ethics.

\item {\bf Broader impacts}
    \item[] Question: Does the paper discuss both positive and negative societal impacts of the work performed?
    \item[] Answer: \answerYes{}
    \item[] Justification: Section~\ref{sec:conclusion} notes that this work targets a clinical communication gap (expectation management before rhinoplasty) and explicitly requires prospective clinical evaluation before any clinical use. The paper discloses failure modes that could mislead patients (GT-below-BL finding, small $N$, no cosmetic panel study), the non-commercial FLUX.1-dev license constraint, the Monk Skin Tone coverage gap, and the dual-use risk of photorealistic face simulation. No positive framing is offered without the corresponding limitation.
    \item[] Guidelines: \answerYes{} if the paper discusses both positive and negative societal impacts.

\item {\bf Safeguards}
    \item[] Question: Does the paper describe safeguards for responsible release?
    \item[] Answer: \answerYes{}
    \item[] Justification: Procedural safeguards: the released code does not include model weights; HDA images are not redistributed; external validation images are not bundled; the anonymous code repository contains only evaluation harnesses and preset configurations. Architectural safeguards: the hard-mask composite restricts generation to the surgical region; the identity gate ($\geq 0.65$ ArcFace) rejects outputs with identity drift; the deterministic TPS fallback (M5) activates when all diffusion candidates fail. The paper includes a prominent disclaimer that outputs are simulation only and not for clinical use, and notes the non-commercial FLUX.1-dev license.
    \item[] Guidelines: \answerYes{} if safeguards are described for assets that might be misused.

\item {\bf Licenses for existing assets}
    \item[] Question: Are existing assets properly credited with licenses?
    \item[] Answer: \answerYes{}
    \item[] Justification: Key assets are cited with license information: FLUX.1-dev and FLUX.1-Fill-dev operate under the FLUX.1-dev non-commercial license; the HDA database is used under its published research-use terms~\citep{rathgeb2020plastic}; InsightFace buffalo\_l is used under its research-use restriction; MediaPipe operates under Apache 2.0~\citep{kartynnik2019mediapipe}; Depth Anything V2 operates under Apache 2.0~\citep{yang2024depth}; LPIPS uses the MIT license~\citep{zhang2018lpips}; ControlNet weights are from jasperai under the FLUX non-commercial license.
    \item[] Guidelines: \answerYes{} if all assets are cited and licensed appropriately.

\item {\bf New assets}
    \item[] Question: Are new assets introduced in the paper well documented?
    \item[] Answer: \answerYes{}
    \item[] Justification: Three new assets are released at the supplementary material (zipped code bundle submitted alongside this paper PDF): (1) the SurgicalScore evaluation package with documented API, usage examples, and unit tests; (2) the \method{} reference pipeline with a README and step-by-step reproduction guide; (3) HDA matched test split manifests/indices for $N{=}27$ blepharoplasty, $N{=}21$ rhinoplasty, and $N{=}9$ rhytidectomy paired cases (case identifiers and per-case metric outputs); images are not redistributed and are obtained from the original authors~\citep{rathgeb2020plastic} under their research-use terms. A data card documents source, split procedure, known issues (Facelift\_08 subject mismatch, filtered at ingest), and intended use. The supplementary README is the canonical documentation.
    \item[] Guidelines: \answerYes{} if new assets are documented.

\item {\bf Crowdsourcing and research with human subjects}
    \item[] Question: Does the paper include full instructions given to participants and screenshots of interfaces (if applicable) as well as procedures for ensuring ethical treatment of participants and for data confidentiality?
    \item[] Answer: \answerNA{}
    \item[] Justification: The paper involves no crowdsourcing and no formal human subjects research. The surgeon review described in Section~\ref{sec:results} is an informal visual check by one expert reviewer of generated output quality on six outputs selected to span SurgicalScore range and visible failure modes; it is not a user study, not a clinical study, and no human subjects were enrolled. The reviewer evaluated model outputs, not surgical decisions affecting any patient. This does not constitute human subjects research under standard definitions.
    \item[] Guidelines: \answerNA{} if the paper does not involve crowdsourcing or research with human subjects.

\item {\bf Institutional review board (IRB) approvals}
    \item[] Question: Does the paper describe the institutional review board (IRB) approvals or other review protocols for each dataset used, following the risk and benefit guidelines of the involved countries?
    \item[] Answer: \answerNA{}
    \item[] Justification: All evaluation data are retrospective. HDA and ASPS are publicly accessible or research-use sources; PCA images are used under authorization from the providing clinical archive and are not redistributed. No new human subjects were recruited, no intervention was performed, and prospective deployment would require institutional review and consent procedures.
    \item[] Guidelines: \answerNA{} if the paper does not involve IRB approvals or other review protocols.

\item {\bf Declaration of LLM usage}
    \item[] Question: Does the paper describe usage of LLMs as an important component of the core methods?
    \item[] Answer: \answerNA{}
    \item[] Justification: LLMs are not a component of the core methods. The pipeline uses diffusion models (FLUX.1-dev), a depth estimator (Depth Anything V2), and discriminative models (ArcFace, MediaPipe). LLMs were used for writing assistance only, which per NeurIPS policy does not require declaration in the checklist.
    \item[] Guidelines: \answerNA{} if the paper does not involve LLMs as a core method component.

\end{enumerate}
